\documentclass[12pt]{article}
\usepackage[margin=0.9in]{geometry}
\usepackage{xcolor}
\usepackage{parskip}
\usepackage{enumitem}
\usepackage{titlesec}
\usepackage{fancyhdr}
\usepackage{amsmath,amssymb}
\usepackage[pdfusetitle]{hyperref}
\hypersetup{colorlinks=true, linkcolor=blue!60!black, urlcolor=blue!60!black}

% ===== colors =====
\definecolor{slidebar}{RGB}{30, 70, 130}
\definecolor{action}{RGB}{170, 40, 40}
\definecolor{onscreen}{RGB}{80, 80, 80}
\definecolor{pauseclr}{RGB}{120, 120, 120}
\definecolor{transition}{RGB}{20, 110, 50}

% ===== header =====
\pagestyle{fancy}
\fancyhf{}
\renewcommand{\headrulewidth}{0pt}
\fancyfoot[C]{\small\color{pauseclr}\thepage}

% ===== slide heading =====
\newcommand{\slidehd}[3]{%
  \clearpage
  \noindent\colorbox{slidebar}{\color{white}\parbox{\dimexpr\linewidth-2\fboxsep}{%
    \vspace{0.2em}%
    {\Large\bfseries SLIDE #1 \quad $\cdot$ \quad #2}\\[0.2em]%
    {\small timing: #3}%
    \vspace{0.2em}%
  }}\par\vspace{0.8em}%
}

% ===== on-screen cue =====
\newcommand{\onscreen}[1]{%
  \noindent{\color{onscreen}\footnotesize\textsc{on screen:} \textit{#1}}\par\vspace{0.6em}
}

% ===== verbatim speech =====
\newenvironment{say}{%
  \par\medskip%
  \begin{quote}%
  \setlength{\parskip}{0.7em}%
  \large\itshape%
}{%
  \end{quote}%
  \medskip%
}

% ===== stage directions =====
\newcommand{\act}[1]{%
  \par\smallskip\noindent{\color{action}\textbf{$\rightarrow$\; #1}}\par\smallskip%
}
\newcommand{\pauseact}{%
  \par\smallskip\noindent{\color{pauseclr}\textit{(pause)}}\par\smallskip%
}
\newcommand{\transact}[1]{%
  \par\medskip\noindent{\color{transition}\textbf{$\Rightarrow$\; TRANSITION:} #1}\par%
}
\newcommand{\handoff}[1]{%
  \par\medskip\noindent{\color{transition}\textbf{$\Rightarrow$\; HAND OFF:} #1}\par%
}

% ===== metadata =====
\title{Speaker Script --- Suleyman's Sections}
\author{CAS 760 Development Project 2 \\ A Simple File System in Alonzo}
\date{Presentation day}

\begin{document}

% ===== cover page =====
\begin{titlepage}
\centering
\vspace*{3cm}

{\Huge\bfseries Speaker Script}\par
\vspace{0.5em}
{\Large Suleyman's Sections}\par
\vspace{2em}
{\large CAS 760 Development Project 2}\par
{\large A Simple File System in Alonzo}\par
\vspace{4em}

\begin{minipage}{0.75\textwidth}
\textbf{Target timing:} $\sim$7 minutes total, $\sim$60 seconds per slide.\\[0.6em]
\textbf{My slides, in order (deck slides 8--14):}
\begin{enumerate}[leftmargin=2em,itemsep=0.2em]
  \item Theory for Files --- \textsc{FILE}
  \item Theory Extension for Secured Files --- \textsc{SECURED-FILE}
  \item Theory Extension for Directories --- \textsc{DIRECTORY}
  \item Directories as an Inductive Type
  \item Directory Development --- \textsc{DIR-DEV} (both \texttt{contains} and \texttt{size} defined here)
  \item Explanation for Directory Development
  \item Example of Directory Evaluation (handoff to Hamza)
\end{enumerate}

\vspace{1.5em}
\textbf{How to read this:}
\begin{itemize}[leftmargin=1.5em,itemsep=0.2em]
  \item The \textit{italic indented text} is what I say \textbf{verbatim}.
  \item {\color{action}\textbf{$\rightarrow$ Red arrows}} are stage directions --- where to point, gesture, or pause.
  \item {\color{pauseclr}\textit{(pause)}} means take a breath, about one second.
  \item {\color{transition}\textbf{$\Rightarrow$ Green TRANSITION}} means advance to the next slide.
  \item Each slide starts on a fresh page --- flip as you go.
\end{itemize}
\end{minipage}

\vfill
{\footnotesize\color{pauseclr} Read this aloud at least twice before the talk. Time yourself.}
\end{titlepage}

% ============================================================
\slidehd{1}{Files --- Theory \textsc{FILE}}{$\sim$70 sec}

\onscreen{Base types $N, C, F$. Constants $\mathrm{mkfile}, \mathrm{fname}, \mathrm{fcontent}$ with types. Seven axioms (1--7) listed.}

\begin{say}
``Thanks Hamza. Now I'll walk through the Files theory.''
\end{say}
\act{Brief nod toward Hamza. Face the audience.}

\begin{say}
``The idea here is straightforward: a file is a named container for content. Conceptually, a file is just a pair --- a name and some content --- but we want to formalize that abstractly, without committing to what names or contents actually are.''
\end{say}

\begin{say}
``So we introduce three base types ---''
\end{say}
\act{Point to the line ``$N,\; C,\; F$'' on the slide.}
\begin{say}
``--- $N$ for names, $C$ for content, and $F$ for the file type itself. We keep them opaque on purpose, so that this theory captures the \textbf{structure} of files rather than any specific representation.''
\end{say}

\pauseact

\begin{say}
``We have one constructor, $\mathrm{mkfile}$ ---''
\end{say}
\act{Point to the $\mathrm{mkfile}$ line.}
\begin{say}
``--- of type $N \to C \to F$, which builds a file from a name and content. And we have two accessors, $\mathrm{fname}$ and $\mathrm{fcontent}$ ---''
\end{say}
\act{Point to the $\mathrm{fname}$ and $\mathrm{fcontent}$ lines in turn.}
\begin{say}
``--- that pull the name and content back out.''
\end{say}

\begin{say}
``The axioms say what you'd expect. Axioms 1 and 2: $\mathrm{mkfile}$ is total and injective --- every name-content pair gives a valid file, and distinct pairs give distinct files. Axioms 3 and 4 are the two accessor axioms: if you build a file with $\mathrm{mkfile}\; n\; c$, extracting the name gives you $n$ back, and the content gives you $c$. Axioms 5 and 6 say the accessors are themselves total. And Axiom 7 --- surjectivity --- says every file $f$ is equal to $\mathrm{mkfile}$ applied to its own name and content. That closes the loop: $F$ is isomorphic to $N \times C$, but presented abstractly as a theory.''
\end{say}

\transact{advance to Secured Files slide.}

% ============================================================
\slidehd{2}{Secured Files --- Theory Extension \textsc{SECURED-FILE}}{$\sim$75 sec}

\onscreen{``Extends: FILE and PERM''. ``New base types: none ($P$ inherited from PERM)''. New constants $\mathrm{perm}, \mathrm{protect}$. Axioms 8--12.}

\begin{say}
``SECURED-FILE combines files and permissions. This is one of the more interesting design points in the paper.''
\end{say}

\begin{say}
``The intuition is: FILE plus PERM equals SECURED-FILE. Each file gets a permission level attached to it. Formally, we present SECURED-FILE as a \textbf{joint extension} of both FILE and PERM ---''
\end{say}
\act{Point to the ``Extends: FILE and PERM'' line on the slide.}
\begin{say}
``--- this is a pushout-style construction in the little theories method.''
\end{say}

\begin{say}
``What that means concretely: SECURED-FILE has \textbf{no new base types} of its own.''
\end{say}
\act{Point to ``New base types: none'' line.}
\begin{say}
``It inherits $P$ from PERM and inherits $N$, $C$, $F$ from FILE. The two extension arrows in the theory graph --- from PERM to SECURED-FILE, and from FILE to SECURED-FILE --- encode that inheritance formally.''
\end{say}

\pauseact

\begin{say}
``What SECURED-FILE \textbf{does} add is two new constants.''
\end{say}
\act{Point to $\mathrm{perm}$ line.}
\begin{say}
``First, $\mathrm{perm}$, of type $F \to P$, which tells you the permission of a file.''
\end{say}
\act{Point to $\mathrm{protect}$ line.}
\begin{say}
``And second, $\mathrm{protect}$, of type $F \to P \to F$, which takes a file and a permission and gives you back the same file with that permission assigned.''
\end{say}

\begin{say}
``Axioms 8 through 12 govern these. Eight and nine are totality. Ten is the core semantics: applying $\mathrm{protect}$ with permission $p$ gives a file whose permission is $p$. And eleven and twelve say that $\mathrm{protect}$ only changes the permission --- it preserves the name and the content. So $\mathrm{protect}$ is a pure permission-update operation.''
\end{say}

\transact{advance to Directories slide.}

% ============================================================
\slidehd{3}{Directories --- Theory Extension \textsc{DIRECTORY}}{$\sim$70 sec}

\onscreen{``Design choice: flat structure.'' Base type $D$. Constructors $\mathrm{emptydir}, \mathrm{addfile}$. Axioms 13--16.}

\begin{say}
``DIRECTORY is a theory extension over SECURED-FILE. This is where the design choice I want to highlight comes in: we model directories as a \textbf{flat} structure, not as a hierarchical tree.''
\end{say}

\begin{say}
``The reason is that this keeps the inductive type clean.''
\end{say}
\act{Point to the $\mathrm{addfile}$ type signature on the slide.}
\begin{say}
``The constructor $\mathrm{addfile}$ has type $F \to D \to D$ --- it takes a \textbf{file}, not a directory, as its first argument. That gives us a direct analogy with the list type from the textbook: $\mathrm{emptydir}$ plays the role of $\mathrm{nil}$, and $\mathrm{addfile}$ plays the role of $\mathrm{cons}$.''
\end{say}

\begin{say}
``If we wanted a hierarchical model, we'd need a second constructor, $\mathrm{adddir}$ of type $D \to D \to D$ ---''
\end{say}
\act{If whiteboard available: quickly write ``\texttt{adddir : D $\to$ D $\to$ D}'' in the corner. Otherwise just emphasize with hand gesture.}
\begin{say}
``--- which would make this a tree-shaped inductive type. That's a natural extension, and we mention it as future work in the paper, but for this project the flat model is the right scope.''
\end{say}

\pauseact

\begin{say}
``The axioms are standard. Axiom 13 says $\mathrm{addfile}$ is total. Axiom 14 says it's \textbf{injective} --- and we'll see later that this injectivity is what makes the FILE-to-DIR morphism work.''
\end{say}
\act{Slight emphasis on ``injective'' --- this callback matters for Hamza's morphism section.}
\begin{say}
``Axiom 15 is disjointness: $\mathrm{addfile}\; f\; d$ is never equal to $\mathrm{emptydir}$. And Axiom 16 is the structural induction principle, which lets us prove properties over all directories by handling the two constructor cases.''
\end{say}

\transact{advance to Inductive Type verification slide.}

% ============================================================
\slidehd{4}{Directories as an Inductive Type}{$\sim$70 sec}

\onscreen{$K_D = \{\mathrm{emptydir}, (\mathrm{addfile}, U_F)\}$. Interpretation bullets. Conclusion about finite lists of $D_F^M$.}

\begin{say}
``To confirm DIRECTORY is a proper recursive inductive type in Alonzo, we follow the procedure from Chapter 10 of Farmer's textbook.''
\end{say}

\begin{say}
``The constructor set is $K_D$ ---''
\end{say}
\act{Point to the $K_D = \{\ldots\}$ expression on the slide.}
\begin{say}
``--- which contains $\mathrm{emptydir}$ and the pair $(\mathrm{addfile}, U_F)$. $\mathrm{emptydir}$ is a \textbf{first-kind} constructor --- it's nullary. And $(\mathrm{addfile}, U_F)$ is a \textbf{fourth-kind} constructor: the $U_F$ means the first argument ranges over the entire domain of $F$, and because $D$ appears in the domain of $\mathrm{addfile}$'s type, it makes $\mathrm{addfile}$ recursive.''
\end{say}

\pauseact

\begin{say}
``We check the seven obligations from Chapter 10 in the paper. I'll just mention them briefly: the quasitype for the $F$-argument is non-empty because $F$ is a base type; totality and injectivity hold by Axioms 13 and 14; $\mathrm{emptydir}$ is not in the range of $\mathrm{addfile}$ by disjointness; disjoint-ranges is trivial since there's only one non-nullary constructor; and structural induction is Axiom 16. All seven pass.''
\end{say}

\begin{say}
``The consequence, by the Inductive Types Theorem, is that in every standard model $D$ is isomorphic to the set of \textbf{finite lists} over the interpretation of $F$. So directories really do behave like lists of files --- exactly as the constructor signatures suggested.''
\end{say}

\transact{advance to DIR-DEV intro slide.}

% ============================================================
\slidehd{5}{Directory Development --- \textsc{DIR-DEV}}{$\sim$90 sec}

\onscreen{Name: DIR-DEV. Bottom theory: DIR. Def 1 ($\mathrm{contains}$) with Thms 1--2. Def 2 ($\mathrm{size}$) with Thms 3--6. All on one slide.}

\begin{say}
``DIR-DEV is a \textbf{development} over DIRECTORY --- not another theory extension. That matters because $\mathtt{contains}$ and $\mathtt{size}$ are \textbf{derived} operations, definable from $\mathrm{addfile}$ and $\mathrm{emptydir}$. So we \textbf{prove} they exist rather than postulating them as new theory constants. In Project 1 we lost points for making the opposite mistake --- this is the direct correction.''
\end{say}
\act{Brief pause on the "lesson learned" beat --- earns goodwill.}

\begin{say}
``The slide gives both definitions and all six theorems. I'll walk through them.''
\end{say}

\act{Point to Def 1 ($\mathtt{contains}$) on the slide.}
\begin{say}
``Def 1: $\mathtt{contains}$ of type $N \to D \to \mathbb{B}$. Defined as the unique function $g$ satisfying two recursion clauses --- $g\; n\; \mathrm{emptydir}$ is false, and $g\; n\; (\mathrm{addfile}\; f\; d)$ is true if $\mathrm{fname}\; f$ equals $n$, or if $g\; n\; d$ is already true. So $\mathtt{contains}$ either matches the outermost file's name, or recurses into the tail.''
\end{say}

\act{Move pointer to Thms 1--2.}
\begin{say}
``Theorem 1 says $\mathtt{contains}$ is total. Theorem 2 says the empty directory contains no name.''
\end{say}

\pauseact

\act{Point to Def 2 ($\mathtt{size}$) on the slide.}
\begin{say}
``Def 2: $\mathtt{size}$ of type $D \to \mathbb{N}$. Same iota-term pattern --- $g\; \mathrm{emptydir}$ equals zero, and $g\; (\mathrm{addfile}\; f\; d)$ equals $g\; d$ plus one.''
\end{say}

\act{Move pointer down the four theorems 3--6.}
\begin{say}
``Four theorems: Theorem 3, totality. Theorem 4, size of empty is zero. Theorem 5, the recurrence --- $\mathtt{size}\; (\mathrm{addfile}\; f\; d)$ equals $\mathtt{size}\; d$ plus one. And Theorem 6, non-negativity, which follows from the recurrence by structural induction. Both definitions are validated in the appendix.''
\end{say}

\transact{advance to the Explanation slide.}

% ============================================================
\slidehd{6}{Explanation for Directory Development}{$\sim$50 sec}

\onscreen{Plain-language prose summary of $\mathtt{contains}$ and $\mathtt{size}$. Two bullet blocks. Mentions appendix validation.}

\begin{say}
``This slide restates the same two definitions in plain language for intuition.''
\end{say}

\act{Point to the \texttt{Contains} block on the slide.}
\begin{say}
``$\mathtt{contains}\; n\; d$ asks: is there a file named $n$ somewhere in directory $d$? The empty directory contains nothing; and $\mathrm{addfile}\; f\; d$ contains $n$ exactly when either $f$'s own name is $n$, or $d$ already contained $n$.''
\end{say}

\act{Point to the \texttt{Size} block.}
\begin{say}
``$\mathtt{size}\; d$ counts file entries: zero for $\mathrm{emptydir}$, and one more than $\mathtt{size}\; d$ whenever we prepend a file with $\mathrm{addfile}$.''
\end{say}

\pauseact

\begin{say}
``Both definitions are iota-terms --- definite descriptions. The appendix discharges existence and uniqueness for each, which is what makes them legal development-level abbreviations.''
\end{say}

\transact{advance to the concrete example slide.}

% ============================================================
\slidehd{7}{Example of Directory Evaluation}{$\sim$70 sec}

\onscreen{Two files $f_1, f_2$ built via $\mathrm{protect}(\mathrm{mkfile}\; n\; c)\; p$. Two-entry directory $d$. Values of $\mathtt{contains}$ at $n_1, n_2, n_3$ and $\mathtt{size}\; d = 2$.}

\begin{say}
``To close out, a concrete evaluation.''
\end{say}

\act{Point to $f_1$ and $f_2$ definitions on the slide.}
\begin{say}
``Let $f_1$ be $\mathrm{protect}$ applied to $\mathrm{mkfile}\; n_1\; c_1$ with \textbf{read} permission. Let $f_2$ be $\mathrm{protect}$ applied to $\mathrm{mkfile}\; n_2\; c_2$ with \textbf{write} permission. Let $d$ be $\mathrm{addfile}\; f_1\; (\mathrm{addfile}\; f_2\; \mathrm{emptydir})$ --- a two-entry directory.''
\end{say}

\act{Point to each $\mathtt{contains}$ equation.}
\begin{say}
``$\mathtt{contains}\; n_1\; d$ is \textbf{true} --- $f_1$'s name matches at the outer layer. $\mathtt{contains}\; n_2\; d$ is also \textbf{true} --- we recurse into the tail and find $f_2$. $\mathtt{contains}\; n_3\; d$ is \textbf{false} --- no match at either layer.''
\end{say}

\act{Point to the $\mathtt{size}$ line.}
\begin{say}
``And $\mathtt{size}\; d$ equals \textbf{two} --- each $\mathrm{addfile}$ contributes one, $\mathrm{emptydir}$ contributes zero. So the two-entry directory has size two, confirming Theorem 5.''
\end{say}

\begin{say}
``That completes my sections of the development. Hamza will now take over with the theory morphism.''
\end{say}

\handoff{look to Hamza, step slightly aside, smile.}

% ============================================================
\clearpage
\noindent\colorbox{slidebar}{\color{white}\parbox{\dimexpr\linewidth-2\fboxsep}{%
  \vspace{0.2em}%
  {\Large\bfseries Q\&A PREP $\cdot$ Likely Questions}\\[0.2em]%
  {\small rehearse these --- one of them almost certainly comes up}%
  \vspace{0.2em}%
}}\par\vspace{0.8em}%

\textbf{\color{slidebar}Q1. Why flat directories and not hierarchical?}

\begin{say}
``Scope --- the project name says \textbf{simple} file system. Flat keeps the inductive type within the standard fourth-kind constructor pattern of Chapter 10, preserves the direct list analogy, and makes the FILE-to-DIR morphism clean because single-entry directories give a canonical target for the embedding. Hierarchy would require a second recursive constructor $\mathrm{adddir} : D \to D \to D$, which breaks the clean list analogy and makes the morphism extraction ambiguous. We note hierarchy as future work in the paper.''
\end{say}

\vspace{0.8em}
\textbf{\color{slidebar}Q2. Why are \texttt{contains} and \texttt{size} in a development instead of extending DIRECTORY?}

\begin{say}
``Because they're \textbf{derived} operations --- they're definable from the primitives $\mathrm{addfile}$ and $\mathrm{emptydir}$, so postulating them as theory constants would add redundant axioms. The development framework lets us define them once, prove they exist uniquely, and then use them. In Project 1 we put derived operations into theory extensions and lost points for it --- this is the direct correction.''
\end{say}

\vspace{0.8em}
\textbf{\color{slidebar}Q3. How did you verify the DIRECTORY obligations?}

\begin{say}
``All seven obligations from Chapter 10 of Farmer's textbook --- non-empty quasitypes, totality on domain, injectivity, disjointness of ranges, structural induction. Each one discharges directly from one of our four axioms, Axioms 13 through 16, with the non-emptiness obligation for the $F$-argument holding because $F$ is a base type and every base type is inhabited in any Alonzo frame. The full check is in the paper.''
\end{say}

\vspace{0.8em}
\textbf{\color{slidebar}Q4. What if two files in a directory have the same name?}

\begin{say}
``They're allowed. The inductive type doesn't enforce uniqueness --- you can have two $\mathrm{addfile}$ applications with files sharing a name. $\mathtt{contains}$ would just return true if \textbf{any} match is found. A future extension could add a uniqueness invariant, but that's a development-layer property, not a theory-layer one.''
\end{say}

\vspace{0.8em}
\textbf{\color{slidebar}Q5. Why use $K_D$ for the constructor set?}

\begin{say}
``Standard inductive-type convention --- $K$ for \textit{Konstruktor}, subscripted by the type. The subscript makes it unambiguous which type's constructors we mean, and it avoids clashing with the base type $C$ of FILE.''
\end{say}

\vspace{0.8em}
\textbf{\color{slidebar}Q6. What does ``pushout-style construction'' mean?}

\begin{say}
``It's a way of combining two theories that share no non-trivial sub-theory. SECURED-FILE takes FILE and PERM as inputs and produces a new theory whose language is the union of both, augmented with the new constants $\mathrm{perm}$ and $\mathrm{protect}$. The category-theoretic term is a pushout over the empty theory. In Alonzo's little theories method, we encode it by having SECURED-FILE extend both parents simultaneously, which is what the two arrows in our theory graph represent.''
\end{say}

\end{document}
